\section{Introduction}
\label{sec:intro}
\subsection{Background}

In the present paper we deal with summation 
identities for spin Hall--Little\-wood symmetric
rational functions. These functions
arise
as 
partition functions
of square lattice integrable
vertex models
related to the quantum group $U_q(\widehat{\mathfrak{sl}_2})$. This description originally appeared in \cite{Bor17,BP18}.

The spin Hall--Littlewood functions
also can be identified with Bethe Ansatz eigenfunctions 
of the higher spin six vertex model on $\mathbb{Z}$, \cf \cite[Ch. VII]{KBI93}.
They also appear as eigenfunctions of certain stochastic particle systems
\cite{Pov13,BCPS15,CP16}. Following
\cite{Bor17,BP18} and subsequent works, we treat spin Hall--Littlewood functions and their relatives from the point of
view of the theory of symmetric functions.
A classical reference on the theory of symmetric functions is the book
\cite{Mac95} where Schur, Hall--Littlewood, and Macdonald symmetric
polynomials
and symmetric functions are developed and various identities for them are formulated or proved.


One of the common features 
for most
families of symmetric polynomials
is a \emph{Littlewood type summation identity}. For example, the Schur symmetric polynomials $s_{\lambda}$ satisfy the following Littlewood identity:
\begin{equation*}
\sum_{\lambda'\ \text{even}} s_{\lambda}(u_1, \dots, u_m)
=
\prod_{1 \leq i<j\leq m}
(1 - u_i u_j)^{-1}.
\end{equation*}
Here the summation is over all signatures $\lambda=(\lambda_1\ge \ldots\ge \lambda_m\ge0 )$ such that all parts of its conjugate $\lambda'$ are even, or equivalently, all part-multiplicities of $\lambda$ are even.
For a comprehensive study of Littlewood identities for
Hall–Littlewood polynomials, we refer the reader to \cite[Ch. III]{Mac95} and to \cite{War06,RW21} for recent developments concerning boxed Littlewood formulae for
Macdonald polynomials.

Moreover, Littlewood identities are important for integrable probability: they appear as a key tool for studying half-space integrable models related to the corresponding half-space Macdonald processes, see \cite{BBC20,BBCW18,BZ19}. Also, these types of identities for stable spin Hall--Littlewood polynomials, which are specializations of our functions, were applied in \cite{CD21} for introducing the half-space Yang-Baxter random field and studying related dynamic systems.

We study \emph{refinements} of Littlewood type identities,
which are derived by inserting an extra factor
into each term of the summation in the left-hand side.
The expression for the right-hand side, in turn, also gets more complicated: it becomes a Pfaffian. Earlier, a number of Pfaffian formulas for partition functions of the six vertex model were obtained by Kuperberg in \cite{Kup02}. We follow a method for proving refined (Cauchy and Littlewood type) identities introduced in \cite{WZJ16}, which is based on the Yang--Baxter equation.

One of the applications of refined Cauchy identities for Macdonald polynomials is a possibility to compute the expectations of observables for Macdonald
measures. Namely, they can be expressed as a certain determinantal formula independent of the parameter $q$.
This result goes back to 
\cite{KN99}, see also \cite{War08,Bor18,Pet21}.
It would be nice to see if this $q$-independence extends to the Pfaffian case. Moreover, it would be interesting to employ our result for the analysis of half-space models of integrable probability as in \cite{BBC20,BBCW18,BZ19}, but this application is outside of the scope of the present work.

\subsection{Refined Littlewood identity for spin Hall--Littlewood functions}

One of possible ways to define the fully inhomogeneous 
spin Hall--Littlewood 
symmetric rational functions is the following
symmetrization form introduced in \cite{BP18}:
\begin{equation*}
\begin{split}
F_\lambda(u_1,\ldots,u_N )
&=
\sum_{\sigma\in S_N}
\sigma
\Biggl( 
\prod_{1\le i<j\le N}\frac{u_i-t u_j}{u_i-u_j}\,
\prod_{i=1}^{N}
\biggl(
\frac{1-t}{1-s_{\lambda_i}u_i}
\prod_{j=0}^{\lambda_i-1}\frac{ u_i-s_j}{1- s_j u_i}
\biggr)
\Biggr),
\end{split}
\end{equation*}
where 
$\lambda=(\lambda_1\ge\ldots\ge \lambda_N\ge0 )$ is a signature, that is, a sequence of weakly decreasing nonnegative integers.
Here $\sigma\in S_N$
acts by permuting the variables $u_i$'s.
The function $F_\lambda$ depends on the ``quantum parameter''
$t \in (0, 1)$, the variables $u_j$ and the inhomogeneities $s_x$, where $x\in \mathbb{Z}_{\ge0}$.
By setting $s_x=0$ for all $x$, we obtain the reduction to the case of usual
Hall--Littlewood symmetric polynomials.

Our main result is a generalization of the refined Littlewood identity \eqref{class_id} to the case of the spin Hall--Littlewood functions. 

To formulate the result we need some notation given below. Namely, $m_0(\lambda)$ is the number of parts in signature $\lambda$ equal to zero, and 
$(a;t)_k=(1-a)(1-at)\ldots(1-at^{k-1})$ is the $t$-Pochhammer symbol.
\begin{theorem} \label{lit_id} 
	Let $\gamma\ne0$ be an arbitrary complex number and let variables $u_1, \dots, u_{2n}$ satisfy restrictions \eqref{eq:admissible_u} below which are needed for some convergence conditions.
	Then spin Hall--Littlewood symmetric rational functions satisfy the following refined Littlewood identity:
	\begin{equation} \label{littlewood_id}
	\begin{split}
	\sum_{\lambda: m_{i}(\lambda) \in 2 \mathbb{Z}_{\geq 0}} &\frac{1}{( t; t)_{m_{0}(\lambda)}} \prod_{j = 1}^{m_{0}(\lambda) / 2} (1 - s_{0}^2 \gamma^{-1} t^{2j - 2})(1 - \gamma t^{2j-1})  \prod_{j=1}^{2n} (1 - s_0 u_j)		\\&\times
	\prod_{i = 1}^{\infty} \prod_{j = 1}^{m_{i}(\lambda) / 2} \frac{1 - s_{i}^2 t^{2j - 2}}{1 - t^{2j}} F_{\lambda} (u_1, \dots, u_{2n})
	\; \;\;
	= 
	\prod_{1 \leq i<j \leq 2n}
	\left(
	\frac{1-t u_i u_j}{u_i - u_j}
	\right)
	\\\times
	\pf_{1\leq i < j \leq 2n}&
	\left[
	\frac{ (u_i - u_j)((1 - t)(1 - s_0 u_i)(1- s_0 u_j) + (1 - \gamma)(t-s_{0}^{2}\gamma^{-1})(1-u_i u_j))}
	{(1-u_i u_j) (1-t u_i u_j)}
	\right].
	\end{split}
	\end{equation}
\end{theorem}


\subsection{Sketch of proof}
Our approach follows the work of M. Wheeler and P. Zinn-Justin \cite{WZJ16} and the work of L. Petrov \cite{Pet21}. Namely, we represent spin Hall--Littlewood rational functions as certain partition functions, using the integrable model of deformed bosons. 
Then we consider a partition function that can be identified with some weighted sum of spin Hall--Littlewood functions. After that we use the Yang--Baxter equation to equate our partition function with a partition function of the six vertex model with finitely many vertices. This description allows us to prove certain properties of the partition function and to present a particular function (in our case it is some certain Pfaffian) with the same properties.
Finally, we use Lagrange interpolation to verify that our properties determine the function uniquely. This technique goes back to Izergin and Korepin \cite{Ize87,KBI93}.

Note that in our case we deal with a rather complicated Pfaffian. The novelty of our approach is to simplify one of the steps of the proof. Namely, instead of computing our Pfaffian of order $2n$ in some nontrivial point, which is quite difficult, we specialize our partition function in another point. This makes our argument combinatorial, while the nontrivial pfaffian itself is computed as a by-product of the proof.


\subsection{Notation}
\label{sec:not}

Let us introduce some notation.

Each signature $\lambda=(\lambda_1\ge \lambda_2\ge \ldots\ge \lambda_N\ge0 )$ can be written in multiplicative form as $\lambda=(0^{m_0} 1^{m_1} 2^{m_2} \dots)$, where $m_i$ is the multiplicity of $i$ in $\lambda$. Throughout the paper we will use this notation.

We often deal with tensor products of the same space, so we use upper indices to point out in which component a certain operator acts. For example, if $\omega$ is a  $4\times4$ matrix and we have a $2^n$-dimensional tensor power of $n$ $2$-dimensional spaces, then $\omega^{(i, {i+1})}= \mathbbm{1}^{\otimes ({i - 1})} \otimes \omega \otimes \mathbbm{1}^{\otimes ({n - i - 1})}$ where $\mathbbm{1}$  is a $2\times2$ identity matrix.


\subsection{Organization of the paper} In Section~\ref{sec:vertex_and_sHL} we recall the basic notation,
definitions and properties of the spin Hall--Littlewood rational symmetric functions and the integrable model related to them.
In Section~\ref{sec:sHL_proof} we prove the refined Littlewood identity 
for the spin Hall--Littlewood functions.
Finally, in Section~\ref{sec: spec}
we discuss the reduction of the result to the classical family of Hall-Littlewood symmetric functions and write the non-refined case of our identity.


\section{Higher spin six vertex model weights and spin Hall--Littlewood functions} \label{sec:vertex_and_sHL}

In this section we introduce a certain model with higher spin six vertex weights and explain how to build spin Hall--Littlewood functions in terms of this model.

\subsection{Definition of the model}
Consider an infinite dimensional vector space $V$:
\begin{align*}
\label{space}
V 
=
{\rm Span} 
\left\{
\ket{m_0}_0
\otimes
\ket{m_1}_1 
\otimes
\ket{m_2}_2 
\otimes
\cdots
\right\},
\qquad
m_i \in \mathbb{Z}\ \forall\ i \geq 1,
\end{align*}
where only finitely many of the $m_i$ are nonzero.
In the case where all  $m_i$ are nonnegative, it is convenient to regard the corresponding basis vectors as a particle system on $\mathbb{Z}_{\geq 0}$ in which any number of particles can occupy sites located at each of the positive integers. Namely, each $m_i$ represents just the number of particles at site $i$.
Note that
if $m_i$ are obtained as multiplicities of some signature $\lambda$, then obviously we have $m_i \in \mathbb{Z}_{\geq 0}$  and in this case we denote the corresponding state by $|\lambda\rangle \in V$ or $\langle \lambda | \in V^*$. However, for our purposes it makes sense to work with negative integers, too.

Also, we set up a two-dimensional auxiliary vector space $W = \mathbb{C}^2$ and its basis denoted by
$
|0\rangle $ and $ |1\rangle
$.
Then the higher spin six vertex model weights $w_{u,s_i}(i_1,j_1;i_2,j_2)$ can be seen as the matrix elements of the operators $L_{u, s_i}$ acting in $W \otimes V_i$ where by $V_i$ we denote the $i^{th}$ factor of $V$. Namely,
the weight $w_{u,s}(i_1,j_1;i_2,j_2)$ is defined as $ \langle j_1| \otimes \langle i_1| \; L_{u, s} \; |j_2 \rangle \otimes |i_2\rangle $, where $i_1,i_2\in \mathbb{Z}_{\ge0}$
and $j_1,j_2\in \left\{ 0,1 \right\}$.
Graphically, we can represent this operator as in Figure~\ref{fig:w_weights}. Exact expressions for the weights $w_{u,s}$, except zero ones, are also given in Figure~\ref{fig:w_weights}. Apart from $u$ and $s$, they depend on some fixed parameter $t \in (0, 1)$.
\begin{figure}[htpb]
	\centering
	\includegraphics[scale=0.88]{figures/Table_Weights_HS6VM.pdf}
	\caption{Vertex weights $w_{u,s}(i_1,j_1;i_2,j_2)$.  Here $i_1,i_2\in \mathbb{Z}_{\ge0}$
		and $j_1,j_2\in \left\{ 0,1 \right\}$.}
	\label{fig:w_weights}
\end{figure}

Likewise, let us define operators $L^*_{v, s}$ and the corresponding weights as follows:
\begin{equation}
\label{eq:w_wstar_relation}
w_{v,s}^*(i_1,j_1;i_2,j_2)= \langle j_1| \otimes \langle i_1| \; L^*_{v, s} \; |j_2 \rangle \otimes |i_2\rangle=\frac{(s^2;t)_{i_1}(t;t)_{i_2}}{(s^2;t)_{i_2}(t;t)_{i_1}}\,w_{v,s}(i_2,j_1;i_1,j_2).
\end{equation}
Exact expressions for the nonzero weights $w_{v,s}^*$, together with graphical illustrations, can be found in Figure~\ref{fig:w_star_weights}.

\begin{figure}[htpb]
	\centering
	\includegraphics[scale=0.9]{figures/table_dual_weights.pdf}
	\caption{Vertex weights $w_{v,s}^*(i_1,j_1;i_2,j_2)$. 	Here $i_1,i_2\in \mathbb{Z}_{\ge0}$
		and $j_1,j_2\in \left\{ 0,1 \right\}$.}
	\label{fig:w_star_weights}
\end{figure}

Impose the following restrictions on the variables $u_1, \dots, u_N$:
\begin{equation}
\label{eq:admissible_u}
\left|
\frac{ u_i-s_x}{1-s_x u_i}
\right|\le 1-\varepsilon<1
\qquad \textnormal{for all $i$ and all $x=0,1,2,\ldots $.}
\end{equation}
Here $\varepsilon$ is some fixed positive real number.

Define the following transfer matrices acting on $W \otimes V$:
\begin{align*}
T(u)
&=
\prod_{i=0}^{\infty}
L_{u, s_i} =
\begin{pmatrix}
T_{-}(x) & 0
\medskip
\\
T_{+}(x) & 0
\end{pmatrix}
\in
{\rm End}(W \otimes V_0 \otimes V_1 \otimes \cdots),
\\
T^{*}(u)
&=
\prod_{i=0}^{\infty}
L^*_{u, s_i} =
\begin{pmatrix}
T^{*}_{-}(x) & 0
\medskip
\\
T^{*}_{+}(x) & 0
\end{pmatrix}
\in
{\rm End}(W \otimes V_0 \otimes V_1 \otimes \cdots).
\end{align*}
See Figure~\ref{fig:T} for an illustration.
\begin{figure}[htpb]
	\centering
	\includegraphics[scale=0.9]{figures/fig_T_operators.pdf}
	\caption{Graphical representation of $T$ and $T^{*}$ operators.}
	\label{fig:T}
\end{figure}
\begin{rema}
	Since we require the convergence condition \eqref{eq:admissible_u}, it follows that operators $T_{+}$ and $T^{*}_{+}$ have the vanishing property. Namely, any path that goes endlessly to the right has weight equal to zero, so we can forbid such paths. This means that we do not actually need to write $0$ on the right boundary in Figure~\ref{fig:T}.
\end{rema}

Let us introduce the $R$-matrix of the six vertex model:
\begin{align*}
\label{Rmat}
R_{z}
=
\begin{pmatrix}
1 & 0 & 0 &  \frac{(1-t)z}{1-z} \\
0 & \frac{1-t z}{1-z} & 0 & 0 \\
0 & 0 & \frac{1-t z}{1-z} & 0 \\
\frac{1-t}{1-z} & 0 & 0 & t
\end{pmatrix}
\in 
{\rm End}(W \otimes W).
\end{align*}

Graphically, we denote the action of this operator by cross vertices with the weights given below:
\begin{figure}[htpb]
	\centering
	\includegraphics[scale=0.85]{figures/fig_table_r_big.pdf}
	\caption{Cross vertex weights $R_z$. Here we have $i_1,j_1,i_2,j_2\in \left\{ 0,1 \right\}$.}
	\label{fig:R_weights}
\end{figure}

\begin{prop}[Yang--Baxter equation]
	\label{prop:YBE_w_wstar}
	For $i_1,i_2,j_1,j_2\in  \left\{ 0,1 \right\}$ and 
	$i_3,j_3\in \mathbb{Z}_{\ge0}$ we have
	\begin{equation}
	\label{eq:YBE_w_wstar}
	\begin{split}
	&\sum_{k_1,k_2,k_3}R
	_{u v}(i_2, i_1; k_2, k_1)\,
	w^*_{v,s} (i_3, k_1; k_3, j_1)\,
	w_{u,s}(k_3,k_2; j_3,j_2) \\
	&= 
	\sum_{k'_1,k'_2,k'_3} w^*_{v,s} (k'_3, i_1; j_3, k'_1)\,
	w_{u,s}(i_3,i_2; k'_3,k'_2)\,
	R_{u v}(k'_2, k'_1; j_2, j_1),
	\end{split}
	\end{equation}
	or, equivalently,
	\begin{equation}
	\label{RLL}
	R_{uv}^{(12)} L^{*(1)}_{v, s} L^{(2)}_{u, s} = L^{(1)}_{u, s} L^{*(2)}_{v, s} R_{uv}^{(12)}.
	\end{equation}
\end{prop}

\begin{proof}
	The proof is by direct computations, and we omit them.
\end{proof}

\begin{figure}[htpb]
	\centering
	\includegraphics{figures/fig_YBE_2_colors.pdf}
	\caption{Graphical illustration of the Yang--Baxter equation \eqref{eq:YBE_w_wstar}.}
	\label{fig:YBE_w_wstar}
\end{figure}

\begin{rema}
	It is important that cross vertex weights in the Yang--Baxter equation \eqref{eq:YBE_w_wstar} do not depend on $s$. This observation allows us to iterate this interchange relation horizontally and get an
	equation for $T$ and $T^*$ operators which is the same as \eqref{RLL} with $L_{u, s}$ and $L^{*}_{v, s}$ replaced by $T(u)$ and $T^{*}(v)$, respectively.
\end{rema}

\subsection{Spin Hall-Littlewood functions}
Now we are able to give a definition of the spin Hall--Littlewood rational functions in terms of our model. Namely, they are given by the following formula:
\begin{align*}
F_{\lambda}(u_1,\dots,u_N)
&=
\bra{0}
T_{+}(u_1)
\dots
T_{+}(u_N)
\ket{\lambda}.
\end{align*}

In other words, we consider the weighted sum over all the up-right paths ensembles in $\mathbb{Z}_{\ge0}\times \left\{ 1,\ldots,N  \right\}$ with the following properties:

\begin{enumerate}[\bf1.\/]

\item Each path comes from the left edge; the path entering in row $i$ reaches the top boundary at the corresponding coordinate $\lambda_i$.

\item No two paths can share the same horizontal line.

\item In the vertex $(x, i) \in \mathbb{Z}_{\ge0}\times \left\{ 1,\ldots,N  \right\}$ we take the weight to be $w_{u_i, s_x}$.

\end{enumerate}
An example of such an ensemble is given in Figure~\ref{fig:F_ex}.
\begin{figure}[htpb]
	\centering
	\includegraphics{figures/fig_F_ex.pdf}
	\caption{An example of a path configuration
		contributing to the partition function $F_\lambda(u_1,u_2,u_3,u_4)$, where
		$\lambda=(5,2,2,0)$.}
	\label{fig:F_ex}
\end{figure}

\subsection{Refinement}
One can define a generalization of this partition function by adding an extra parameter $\alpha \in \mathbb{C}$. Namely, consider a vector space \begin{align*}
\label{alpha-space}
V(\alpha) 
=
{\rm Span} 
\left\{
\ket{m_0+ \alpha}_0
\otimes
\ket{m_1}_1 
\otimes
\ket{m_2}_2 
\otimes
\cdots
\right\},
\qquad
m_i \in \mathbb{Z}\ \forall\ i \geq 1,
\end{align*}
and the corresponding family of partition functions $$F^{\alpha}_\lambda(u_1, \dots, u_N) = \bra{0; \alpha}
T_{+}(u_1)
\dots
T_{+}(u_N)
\ket{\lambda; \alpha},$$
where $\lambda$ is a signature and
\begin{align*}
|\lambda; \alpha\rangle
=
\ket{m_0(\lambda) + \alpha}_0
\otimes
\ket{m_1(\lambda)}_1
\otimes
\ket{m_2(\lambda)}_2
\otimes
\cdots
.
\end{align*}

Note that formally we can no longer interpret the quantity $m_0 + \alpha$ as the number of particles on a vertical edge. However, $m_0$ only appears in the weights as the power for the parameter $t$, which allows us to deal with the states $|\lambda; \alpha\rangle$ in the same combinatorial way as with the states $|\lambda\rangle$ but taking into account certain deformations in zero column weights. Thus it makes sense to consider arbitrary $m_0$, including negative ones, which is important for the next section too.

Let $\gamma = t^{\alpha}$. Since the only difference between $F^{\alpha}_\lambda$ and $F_\lambda$ comes from different zero column weights, it is easy to express one partition function in terms of the other: 
\begin{equation}
\label{F_alpha}
\begin{split}
&F^{\alpha}_\lambda(u_1, \dots, u_N) = \frac{(\gamma t; t)_{m_0}}{( t, t)_{m_0}} \prod_{j=1}^{N} \frac{1 - \gamma s_0 u_j}{1 - s_0 u_j}
\left[ F_{\lambda}(u_1, \dots, u_N)
\Big\vert_{\text{$s_0\to \gamma s_0$}}\right],\\
&F_{\lambda}(u_1, \dots, u_N) = F^{0}_{\lambda}(u_1, \dots, u_N).
\end{split}
\end{equation}

\section{Refined Littlewood identity. Proof.}
\label{sec:sHL_proof}

In this section we prove the refined Littlewood identity (Theorem~\ref{lit_id}).
\subsection{A property of the transfer matrices}
\begin{lemma} \label{lem_T_op}
	Consider the following formal weighted sum of all states with even multiplicities:
	\begin{align*}
	\label{e-vect}
	\ket{\rm e;\alpha}
	=
	\sum_{\substack{
			m_{i}(\lambda) \in 2 \mathbb{Z}_{\geq 0} \\ m_0(\lambda) \in 2 \mathbb{Z} 
	}}
	\
	c_{\lambda}(t; \alpha)
	\ket{\lambda;\alpha},
	\end{align*}
	where the weights are given by
	\begin{align*}
	c_{\lambda}(\alpha,t) 
	= 
	\prod_{i = 1}^{\infty} \prod_{j = 1}^{m_{i}(\lambda) / 2} \frac{1 - s_{i}^2 t^{2j - 2}}{1 - t^{2j}}
	\times
	\left\{
	\begin{array}{cl}
	\displaystyle{
		\prod_{j = 1}^{m_{0}(\lambda) / 2} \frac{1 - s_{0}^2 \gamma t^{2j - 2}}{1 - \gamma t^{2j}}
	},
	&
	m_0(\lambda) \geq 0,
	\\ \\ 
	\displaystyle{
		\prod_{j=1}^{-m_0(\lambda)/2}
		\frac{1 - \gamma t^{-2j+2}}{1 - s_{0}^2 \gamma t^{-2j}}
	},
	&
	m_0(\lambda) \leq 0.
	\end{array}
	\right.
	\end{align*}
	
	Then the transfer matrices $T_{\pm}$ and $T^{*}_{\pm}$ have the following property:
	\begin{equation} \label{eq: lem_T_op}
	T_{+}\ket{\rm e;\alpha} = T^{*}_{+}\ket{\rm e;\alpha}, \quad\quad T_{-}\ket{\rm e;\alpha} = T^{*}_{-}\ket{\rm e;\alpha}.
	\end{equation}
\end{lemma}



\begin{proof}
	Take any signature $\mu$ and the corresponding state
	\begin{align*}
	\langle\mu; \alpha|
	=
	\bra{m_0(\mu) + \alpha}_0
	\otimes
	\bra{m_1(\mu)}_1
	\otimes
	\bra{m_2(\mu)}_2
	\otimes
	\cdots
	\end{align*}
	with $m_0(\mu) \in \mathbb{Z}$ and $m_i(\mu) \in \mathbb{Z}_{\geq 0}$ for all $i \geq 1$.
	Note that there exists a unique $\mu_{+}$ with even multiplicities such that $\bra{\rm \mu;\alpha}T_{+} \ket{\rm \mu_{+};\alpha} \neq 0$ or $\bra{\rm \mu;\alpha}T_{-} \ket{\rm \mu_{+};\alpha} \neq 0$ (which of these is nonzero depends on the parity of the sum over all the multiplicities). Also, denote by $\mu_{-}$ the unique signature with even multiplicities such that $\bra{\rm \mu;\alpha}T^{*}_{+} \ket{\rm \mu_{-};\alpha} \neq 0$ or $\bra{\rm \mu;\alpha}T^{*}_{-} \ket{\rm \mu_{-};\alpha} \neq 0$.
	For example, if $\mu = (6, 4, 4, 3, 2, 2, 0)$, then $\mu_{+} = (6, 6, 4, 4, 2, 2, 0, 0)$ and $\mu_{-} = (4, 4, 3, 3, 2, 2)$ (see Figure~\ref{fig:ex} for an illustration).
	\begin{figure}[htpb]
		\centering
		\includegraphics[scale=0.8]{figures/table_example_1.pdf}
		\caption{An illustration for the definition of $\mu_{+}$ and $\mu_{-}$ when $\mu = (6, 4, 4, 3, 2, 2, 0)$.}
		\label{fig:ex}
	\end{figure}
	
	So, we obtain 
	\begin{align*}
	\bra{\rm \mu;\alpha}T_{\pm} \ket{\rm e;\alpha} = c_{\mu_{+}} \bra{\rm \mu;\alpha}T_{\pm} \ket{\rm \mu_{+};\alpha} , \qquad \bra{\rm \mu;\alpha}T^{*}_{\pm} \ket{\rm e;\alpha} = c_{\mu_{-}} \bra{\rm \mu;\alpha}T^{*}_{\pm} \ket{\rm \mu_{-};\alpha}.
	\end{align*}
	It remains to check that \begin{equation}\label{T_pm_rel}
	c_{\mu_{+}} \bra{\rm \mu;\alpha}T_{\pm} \ket{\rm \mu_{+};\alpha} = c_{\mu_{-}} \bra{\rm \mu;\alpha}T^{*}_{\pm} \ket{\rm \mu_{-};\alpha}.
	\end{equation}
	This equality can be seen from the special case of equation~\eqref{eq:w_wstar_relation} and its analogue for the zero column.
	\begin{figure}[htpb]
		\centering
		\includegraphics{figures/table_conn_weights.pdf}
		\includegraphics{figures/table_conn_weights_a.pdf}
		\caption{The correspondence between vertex weights on both sides of \eqref{T_pm_rel}.}
		\label{fig:pr_w}
	\end{figure}
	Indeed, to get $\bra{\rm \mu;\alpha}T_{\pm} \ket{\rm \mu_{+};\alpha}$ with given $\bra{\rm \mu;\alpha}T^{*}_{\pm} \ket{\rm \mu_{-};\alpha}$ we need to do the replacements as in Figure~\ref{fig:pr_w}. These replacements produce some factor, meanwhile the ratio $c_{\mu_{+}}/c_{\mu_{-}}$ precisely compensates this factor. This concludes the proof.
\end{proof}

Graphically, our statement can be represented as in Figure~\ref{fig:T_e}.
\begin{figure}[htpb]
	\centering
	\includegraphics{figures/fig_T_e.pdf}
	\caption{Graphical representation of equation \eqref{eq: lem_T_op}. Here we do not need to specify the state on the left boundary.}
	\label{fig:T_e}
\end{figure}

\subsection{Setting and transformation of the partition function}

Consider the following partition function which has an additional parameter $\alpha$:
\begin{figure}[H]
	\centering
	\includegraphics{figures/fig_part_P.pdf}
	\caption{The left-hand side of the Littlewood identity expressed graphically as a partition function.}
	\label{fig:part_P}
\end{figure}

From the very definition we have
\begin{equation}
\label{part_alpha}
\begin{split}
& 
\mathcal{P}(u_1, \dots, u_{2n}; t, \alpha) \\
&=\sum_{\lambda: m_{i}(\lambda) \in 2 \mathbb{Z}_{\geq 0}} \prod_{j = 1}^{m_{0}(\lambda) / 2} \frac{1 - s_{0}^2 \gamma t^{2j - 2}}{1 - \gamma t^{2j}} \prod_{i = 1}^{\infty} \prod_{j = 1}^{m_{i}(\lambda) / 2} \frac{1 - s_{i}^2 t^{2j - 2}}{1 - t^{2j}} F^{\alpha}_{\lambda} (u_1, \dots, u_{2n}).
\end{split}
\end{equation}

Since we have up-right path ensembles and the left edge is occupied, it follows that only states with non-negative $m_0$ in $\ket{\rm e;\alpha}$ contribute to our summation.

Let us apply Lemma~\ref{lem_T_op} to the upper row. We get the first equality in Figure~\ref{fig:12steps}.
\begin{figure}[htpb]
	\centering
	\includegraphics[scale=0.78]{figures/fig_1st_2nd_steps.pdf}
	\caption{First step of the transformation of the partition function.}
	\label{fig:12steps}
\end{figure}
The second equality in Figure~\ref{fig:12steps} holds due to the completely frozen cross part on the right, which has weight $1$.

Next, using the Yang--Baxter equation, one can move cross part of the partition function to the left edge. Then we repeat this trick several times (see Figure~\ref{fig:34steps} for an illustration).
\begin{figure}[htpb]
	\centering
	\includegraphics[scale=0.78]{figures/fig_3rd_4th_steps.pdf}
	\caption{Next steps of the transformation of the partition function.}
	\label{fig:34steps}
\end{figure}

After moving all the crosses to the left, the partition function factorizes, and the blue frozen part on the right has weight 1, since it is devoid of paths. Thus, we obtain the partition function as in Figure~\ref{fig:5step}. 
\begin{figure}[htpb]
	\centering
	\includegraphics{figures/fig_5th_step.pdf}	\caption{Expression for the left-hand side of the Littlewood identity as a partition function of the inhomogeneous six vertex model with 
		weights $R_{u_iu_j}$ and decorated boundary conditions.}
	\label{fig:5step}
\end{figure}
The boundary vector $\ket{\rm \alpha_{-}}$ involved there can be expressed explicitly as the following weighted sum:
\begin{equation*}
\ket{\rm \alpha_{-}} = \sum_{k=0}^{\infty} \prod_{j=1}^{k} \frac{1 - \gamma t^{-2j+2}}{1 - s_{0}^2 \gamma t^{-2j}} \ket{ \alpha - 2k}.
\end{equation*}


\subsection{Properties of the partition function}
\begin{lemma}
	\label{prop_lem}
	Let $Z_{2n}$ denote the partition function $\mathcal{P} (u_1, \dots, u_{2n}; t, \alpha)$ as in Figure~\ref{fig:5step} multiplied by $\prod_{1 \leq i < j \leq 2n} (1 - u_{i} u_{j}) \prod_{i = 1}^{2n} (1 - s u_i)$.
	Then $Z_{2n}$ possesses the following properties:
	\begin{enumerate}[\bf1.\/]
		\item $Z_{2n}$ is symmetric in $\{u_1, \dots, u_{2n}\}$.
		
		\item $Z_{2n}$ is a polynomial in $u_{2n}$ of degree $2n - 1$.
		
		\item Setting $u_{2n} = u_{2n - 1}^{-1}$, we have the recursion relation
		$$
		Z_{2n}|_{u_{2n} = u_{2n - 1}^{-1}} = (1 - t) (1 - \gamma s_{0} u_{2n}) (1 - \gamma s_{0} u_{2n-1}) \prod_{j = 1}^{2n-2} (1 - t u_{j} u_{2n}) (1 - t u_{j} u_{2n-1}) Z_{2n - 2}.
		$$
		
		\item Under the specialization $u_{2j - 1} = t$, $u_{2j} = 1 / t^2$ for $1 \leq j \leq n$, we have 
		$$Z_{2n}(t,  1 / t^2, \dots, t,   1 / t^2) = \gamma^{n} (t - 1)^{n^{2}}t^{-2n} (-(t - 1/t)^2 )^{n(n-1)/2}   (1 - s_{0} t^{-2})^n (1 -s_{0} t)^{n}.
		$$
		
		\item For $n=1$ we have
		$$Z_{2}(u_1, u_2) = (1 - t)(1 - \gamma s_{0} u_{1}) (1 - \gamma s_{0} u_{2}) + (1 - \gamma) (t - \gamma s_{0}^2) (1 - u_{1} u_{2}).$$
	\end{enumerate}
\end{lemma}

\begin{proof}
	
	To prove that $Z_{2n}$ is symmetric, let us introduce the vertex weights as in Figure~\ref{fig:r_z_1}.
	\begin{figure}[htpb]
		\centering
		\includegraphics[scale=0.85]{figures/fig_table_r.pdf}
		\label{fig:r_z}
	\end{figure}
	\begin{figure}[htpb]
		\centering
		\includegraphics[scale=0.85]{figures/fig_table_r_1.pdf}
		\caption{The vertex weights $r_z$ and $\bar{r}_z$ involved in the proof of symmetry. Here $i_1,j_1,i_2,j_2\in \left\{ 0,1 \right\}$.}
		\label{fig:r_z_1}
	\end{figure}
	One can check that they satisfy the following Yang--Baxter
	equations and the unitary relation:
	\begin{equation} \label{YB_r}
	\begin{split}
	\sum_{k_1,k_2,k_3}&r
	_{u/v}(i_2, i_1; k_2, k_1)\,
	R_{vw} (k_1, k_3; j_1, i_3)\,
	R_{uw}(k_2, j_3;j_2, k_3) \\
	&= 
	\sum_{k'_1,k'_2,k'_3} R_{vw} (i_1, j_3; k'_1, k'_3)\,
	R_{uw}(i_2, k'_3;k'_2; i_3)
	r_{u/v}(k'_2, k'_1; j_2, j_1).
	\end{split}
	\end{equation}
	\begin{equation} \label{YB_r_bar}
	\begin{split}
	\sum_{k_1,k_2,k_3}
	&R_{uw} (i_3, i_2; k_3, k_2)\,
	R_{vw}(k_3, i_1;j_3, k_1)
	\bar{r}_{v/u}(k_2, k_1; j_2, j_1) \\
	&= 
	\sum_{k'_1,k'_2,k'_3}
	\bar{r}
	_{v/u}(i_2, i_1; k'_2, k'_1)\,
	R_{uw} (k'_3, k'_2; j_3, j_2)\,
	R_{vw}(i_3, k'_1;k'_3, j_1).
	\end{split}
	\end{equation}
	\begin{equation}
	\label{eq:YBE_r_bar}
	\begin{split}
	\sum_{k_1,k_2,k_3}&\bar{r}_{v/u}(i_2, i_1; k_2, k_1)\,
	w^*_{v,s} (i_3, k_1; k_3, j_1)\,
	w^*_{u,s}(k_3,k_2; j_3,j_2) \\
	&= 
	\sum_{k'_1,k'_2,k'_3} w^*_{v,s} (k'_3, i_1; j_3, k'_1)\,
	w^*_{u,s}(i_3,i_2; k'_3,k'_2)\,
	\bar{r}_{v/u}(k'_2, k'_1; j_2, j_1).
	\end{split}
	\end{equation}
	\begin{equation} \label{unit_r}
	\sum_{k_1,k_2,l_1, l_2}r
	_{u/v}(i_2, i_1; k_2, k_1)\,
	R_{uv} (k_1; k_2, l_1, l_2)\,
	\bar{r}_{v/u}(l_2; l_1,j_2, j_1) 
	= R_{uv} (i_2, i_1, j_2, j_1).
	\end{equation}
	
	Graphically, these equations can be viewed as in Figures \ref{fig:ybe_3_col} - \ref{fig:unit}:
	\begin{figure}[htpb]
		\centering
		\begin{minipage}{.5\textwidth}
			\centering
			\includegraphics[scale=0.9]{figures/fig_YBE_3_colors.pdf}
			\caption{Graphical illustration of equation \eqref{YB_r}.}
			\label{fig:ybe_3_col}
		\end{minipage}
		\begin{minipage}{.5\textwidth}
		\centering
\includegraphics[scale=0.9]{figures/fig_YBE_3_colors_d.pdf}
\caption{Graphical illustration of equation \eqref{YB_r_bar}.}
\label{fig:ybe_3_col_d}
		\end{minipage}
	\end{figure}
	\begin{figure}[htpb]
		\centering
		\includegraphics{figures/fig_YBE_r_bar.pdf}
		\caption{Graphical illustration of equation \eqref{eq:YBE_r_bar}.}
		\label{fig:ybe_r_bar}
	\end{figure}
	\begin{figure}[htpb]
		\centering
		\includegraphics{figures/fig_mod_unit_rel.pdf}
		\caption{Graphical illustration of equation \eqref{unit_r}.}
		\label{fig:unit}
	\end{figure}
	
	To see property \textbf{1}, add to the partition function in Figure~\ref{fig:5step} a vertex of weight $r_{u_{i+1}/ u_{i}}$ on the left at the $i$-th position, and a vertex of weight $\bar{r}_{u_{i}/u_{i+1}}$ on the right at the $i$-th position. On the one hand, these operations do not change the partition function at all. On the other hand, we can apply Yang--Baxter
	equations several times and the unitary relation to get the partition function with variables $u_i$ and  $u_{i+1}$ swapped. This concludes the proof of Property \textbf{1}.
	
	For property \textbf{2} we may think of $Z_{2n}$ as the partition function for the model given in Figure~\ref{fig:5step} but with the weights $R_{u_iv_j}$ replaced by $(1 - u_i v_j)R_{u_iv_j}$ and the weights $ w^{*}_{u_j}$ replaced by $(1 - s_0 u_j) w^{*}_{u_j}$. This makes all the weights linear, in particular, it allows to verify that $u_{2n}$ contributes to each part of the summation $2n-1$ times with some coefficients (independent of $u_{2n}$).
	
	For property \textbf{3} let us notice that in this case $(1-u_{2n}u_{2n-1})R_{u_{2n}u_{2n-1}}(1,1;1,1) = 0$, so we should avoid this weight at the beginning. The only remaining option is to choose the weight $(1-u_{2n}u_{2n-1})R_{u_{2n}u_{2n-1}}(1,1;0,0) = 1-t$, which leads to factorization of the partition function. After some computations we get the desired property.
	
	To prove property \textbf{4}, note that the chosen $u_1, \dots, u_{2n}$ satisfy 
	$
	1 - t u_i u_j = 0
	$ for all $i, j$ such that $i+j$  is odd. So, for odd $i+j$ we have $R_{u_i v_j}(0,1,0,1) = R_{u_i v_j}(1,0,1,0) = 0$ and $R_{u_i v_j}(k_1,k_1,k_2,k_2) = (-1)^{k_1 + k_2} t^{k_1}$.
	This observation together with some simple freezing/combinatorial arguments implies that there are $2^{n}$ possible configurations with non-zero weights. Moreover, they are uniquely determined by values on the right edge of the six vertex model. One can compute explicitly the weight of each configuration. For example, it can be done through the following recursion relation:
	\begin{equation*}
	\begin{split}
	&Z_{2n}(t,  1 / t^2, \dots, t,   1 / t^2)
	\\
	&=(1 - t) (t^{2n-1} (1 - s_{0}^2 \gamma/t)(1 - \gamma) + (-t)^{2n-1}(1 - s_0 \gamma/ t^{2})(1 - s_0 \gamma t))  
	\\&\qquad\times\prod_{i<2n-1} (1 - u_i u_{2n}) (1 - u_i u_{2n-1})
	Z_{2n-2} (t,  1 / t^2, \dots, t,   1 / t^2).
	\end{split}
	\end{equation*}
	Here the first and the second summands correspond to the cases where we choose $R_{u_{2n} u_{1}}(1, 1, 1, 1)$ or $R_{u_{2n} u_{1}}(1, 1, 0, 0)$ cross vertex weights, respectively. 
	This recursion immediately gives us the formula:
	\begin{equation*}
	\begin{split}
	&
	Z_{2n}(t,  1 / t^2, \dots, t,   1 / t^2) \\&
	=  \gamma^{n} t^{n(n-2)} (1 - s_{0} t^{-2})^n (1 -s_{0} t)^{n}\prod_{i<j} (1 - u_i u_j)|_{(u_1, ... , u_{2n}) = (t,  1 / t^2, \dots, t,   1 / t^2)}.
	\end{split}
	\end{equation*}
	Finally, property {\bf 5} comes from direct computations.
\end{proof}

\subsection{Explicit formula for the partition function}

\begin{theorem} \label{pf_thm}
	The partition function $\mathcal{P} (u_1, \dots, u_{2n}; t, \alpha)$ can be expressed explicitly as follows:
	\begin{equation}
	\begin{split}
	&
	\label{pf_p}
	\mathcal{P} (u_1, \dots, u_{2n}; t, \alpha) \\&
	= 	\prod_{j = 1}^{2n} \frac{1}{1 - s_0 u_j} \prod_{1 \leq i<j \leq 2n}
	\left(
	\frac{1-t u_i u_j}{u_i - u_j}
	\right)
	\pf_{1\leq i < j \leq 2n}
	\left[
	\frac{Z_{2}(u_i, u_j) (u_i - u_j)}
	{(1-u_i u_j) (1-t u_i u_j)}
	\right].
	\end{split}
	\end{equation}
\end{theorem}
\begin{proof}
	First, one can deduce that the Pfaffian on the right-hand side of \eqref{pf_p} multiplied by the product $\prod_{j = 1}^{2n} (1 - s_0 u_j) \prod_{1 \leq i < j \leq 2n} (1 - u_{i} u_{j}) $ satisfy all the properties~\textbf{1}-\textbf{5} from Lemma~\ref{prop_lem}. Namely, we have property~\textbf{1} because both the Pfaffian and the Vandermonde $\prod_{1 \leq i < j \leq 2n} (u_{i} - u_{j}) $ change the sign under swaps $u_i \leftrightarrow u_{i+1}$. Properties~\textbf{2} and~\textbf{5} are straightforward from the very definition of the Pfaffian. To get property~\textbf{3}, one can multiply the last row and column by $\prod_{1 \leq j < 2n} (1 - u_{j} u_{2n}) $ and the second-to-last row and column by $(1 - u_{2n-1} u_{2n})\prod_{1 \leq j < 2n-1} (1 - u_{j} u_{2n-1}) $. Note that all the elements in this matrix hook vanish except two with indices $n-1$ and $n$. In turn,
	\begin{equation*}
	Z_{2}(u_{2n-1},u_{2n})\big\vert_{u_{2n}=u_{2n-1}^{-1}}= (1-t)
	(1-\gamma s_0u_{2n})(1-\gamma s_0 u_{2n-1}).
	\end{equation*}
	Likewise, one can prove property~\textbf{4}, using the recurrence and the following: \begin{equation*}Z_{2}(t,  1 / t^2) = \gamma(1 - t)(1 - s_{0} t^{-2})(1 -s_{0} t).\end{equation*}
	So, it remains to show that these properties determine a function uniquely. For this purpose one can use Lagrange interpolation as in  \cite{Pet21} and \cite[Appendix B]{WZJ16}.
	
	Namely, we assume that two families of polynomials $f_{2n}(u_1, \dots, u_{2n})$ and $g_{2n}(u_1, \dots, u_{2n})$ satisfy properties \textbf{1}-\textbf{5} and prove that $f_{2n} = g_{2n}$ by induction on $n$. The base case  follows from Property \textbf{5}. To prove the induction step, assume that we proved this statement for $n-1$. Then, let us fix $2n-1$ arbitrary non-zero distinct points $u_1, \dots, u_{2n-1}$. Using the recurrence relation and symmetry, we obtain that $f_{2n}$ and $g_{2n}$ treated as polynomials in $u_{2n}$ coincide in $2n-1$  distinct points $u_{1}^{-1}, \dots, u_{2n-1}^{-1}$. Since their degree is $2n-1$, it follows that $f_{2n} - g_{2n} = c \cdot \prod_{i<j}(1 - u_i u_j)$ where $c$ does not depend on $u_{2n}$. However, because of symmetry it does not depend on $u_1, \dots, u_{2n-1}$ either, which means it is an absolute constant. Finally, as can be seen from property~\textbf{4}, $f_{2n}$ and $g_{2n}$ have the same value at a fixed point, hence~$f_{2n} = g_{2n}$.
	
	This concludes the proof of Theorem~\ref{pf_thm}.
\end{proof}

\begin{coro}
	Under the specialization $u_j = t^{2n - j} / (\gamma s_{0})$, we have 
	\begin{equation} \label{pf_p_sp}
	\begin{split}
	&\pf_{1\leq i < j \leq 2n}
	\left[
	\frac{ (t^{j} - t^{i})(\gamma^2 s_0^2(1 - t)(t^{i} - t^{2n})(t^{j}- t^{2n}) + (1 - \gamma)(t-s_{0}^{2}\gamma)(t^{i+j}\gamma^2 s_0^2-t^{4n})}
	{t^{2i+2j-2n}\gamma s_0(\gamma s_0-t^{4n}/(\gamma s_0)) (\gamma s_0 -t^{4n+1-i-j}/(\gamma s_0))}
	\right]
	\\
	&=\prod_{j = 0}^{2n-1} (1 - t^j\gamma^{-1}) \prod_{0 \leq i<j \leq 2n-1}
	\left(
	\frac{t^j - t^i}{\gamma s_0 -t^{i+j+1} / (\gamma s_0)}
	\right)\\
	&\qquad\times
	\mathcal{P}(t^{2n - 1} / (\gamma s_{0}), t^{2n - 2} / (\gamma s_{0}), \dots, 1 / (\gamma s_{0}); t, \alpha)  \\
	&=(-1)^{n} \gamma^{n} t^{n^{2}} \prod_{0 \leq i<j \leq 2n-1}
	\left(
	\frac{t^j - t^i}{\gamma s_0 -t^{i+j+1} / (\gamma s_0)}
	\right)
	\prod_{j=1}^{n}(1 - s_{0}^{2}\gamma t^{-2j+1})(1 - \gamma^{-1} t^{2j-2})
	.
	\end{split}
	\end{equation}
\end{coro}
\begin{proof}
	Consider the lattice interpretation of $\mathcal{P}(t^{2n - 1} / (\gamma s_{0}), \dots, 1 / (\gamma s_{0}); t, \alpha)$. Indeed, under this specialization the right boundary is fixed, and therefore the whole configuration becomes frozen.
	It is not so easy to verify independently that the Pfaffian in the left-hand side of \eqref{pf_p_sp} factorizes.
\end{proof}

Using \eqref{part_alpha} and \eqref{F_alpha}, the left-hand side of the identity \eqref{pf_p} can be rewritten as the weighted sum of $F_{\lambda}$'s over all signatures $\lambda$ with even multiplicities in the following way:
\begin{equation*}
\begin{split}
\sum_{\lambda: m_{i}(\lambda) \in 2 \mathbb{Z}_{\geq 0}} &\frac{1}{( t; t)_{m_{0}(\lambda)}} \prod_{j = 1}^{m_{0}(\lambda) / 2} (1 - s_{0}^2 \gamma t^{2j - 2})(1 - \gamma t^{2j-1}) \prod_{j=1}^{2n} \frac{1 - \gamma s_0 u_j}{1 - s_0 u_j}
\\
&\times \prod_{i = 1}^{\infty} \prod_{j = 1}^{m_{i}(\lambda) / 2} \frac{1 - s_{i}^2 t^{2j - 2}}{1 - t^{2j}} \left[F_{\lambda} (u_1, \dots, u_{2n})\Big\vert_{\text{$s_0\to \gamma s_0$}}\right].
\end{split}
\end{equation*}
After replacing $\gamma s_0$ by $s_0$, we obtain the desired statement of Theorem~\ref{lit_id}.

\section{Some specializations of the Littlewood identity} \label{sec: spec}

In this section we reduce our result
to classical Hall--Littlewood polynomials and we write a non-refined degeneration of our result.
\subsection{Reduction to the case of classical Hall--Littlewood polynomials}
As was shown in \cite{Pet21}, spin Hall--Littlewood rational functions $F_{\lambda}$ can be reduced to classical Hall--Littlewood polynomials $P_\lambda^{HL}$ in the following way:
\begin{equation}
F_{\lambda}(u_1, \dots, u_N)\Big\vert_{\text{$s_x=0$}} = \prod_{r\ge0}(t;t)_{m_r(\lambda)}\cdot
P_\lambda^{HL}(u_1,\ldots,u_N).
\end{equation}
So, after setting $ s_x = 0$ in equation~\eqref{littlewood_id}, we get
\begin{multline} \label{class_id}
\sum_{\lambda: m_{i}(\lambda) \in 2 \mathbb{Z}_{\geq 0}}  \prod_{j = 1}^{m_{0}(\lambda) / 2}(1 - \gamma t^{2j-1})\prod_{i = 1}^{\infty} \prod_{j = 1}^{m_{i}(\lambda) / 2} (1 - t^{2j-1}) P_{\lambda}^{HL} (u_1,\ldots,u_{2n})
\\
=
\prod_{1 \leq i<j \leq 2n}
\left(
\frac{1-t u_i u_j}{u_i - u_j}
\right) 
\pf_{1\leq i < j \leq 2n}
\left[
\frac{ (u_i - u_j)((1 - \gamma t+ (\gamma-1)tu_i u_j)}
{(1-u_i u_j) (1-t u_i u_j)}
\right],
\end{multline}
which coincides with the Littlewood identity proved in \cite[Theorem 5]{WZJ16}.

Note that the classical (non-refined) Littlewood identity can be seen as a special case of \eqref{class_id} at $\gamma = 0$, even though we exclude this point to write down our main and most general Theorem~\ref{lit_id}. Indeed, under this specialization the Pfaffian on the right-hand side is simplified to a product 
$\prod_{1 \leq i<j \leq 2n} \frac{ u_i - u_j}
{1-u_i u_j}$ which gives us the following:

\begin{equation*}
\sum_{\lambda: m_{i}(\lambda) \in 2 \mathbb{Z}_{\geq 0}}  \prod_{i = 1}^{\infty} \prod_{j = 1}^{m_{i}(\lambda) / 2} (1 - t^{2j-1}) P_{\lambda}^{HL} (u_1,\ldots,u_{2n})
=
\prod_{1 \leq i<j \leq 2n}
\frac{1-t u_i u_j}{1 - u_i u_j} .
\end{equation*}

\subsection{Reduction to the unrefined case}
To get unrefined identity, we set $\gamma = 1$ and obtain the following formula:
\begin{multline} \label{unref_id}
\sum_{\lambda: m_{i}(\lambda) \in 2 \mathbb{Z}_{\geq 0}}   \prod_{i = 0}^{\infty} \prod_{j = 1}^{m_{i}(\lambda) / 2} \frac{1 - s_{i}^2 t^{2j - 2}}{1 - t^{2j}} F_{\lambda} (u_1, \dots, u_{2n})
\\
=
\prod_{1 \leq i<j \leq 2n}
\left(
\frac{1-t u_i u_j}{u_i - u_j}
\right) 
\pf_{1\leq i < j \leq 2n}
\left[
\frac{ (u_i - u_j)(1 - t) }
{(1-u_i u_j) (1-t u_i u_j)}
\right].
\end{multline}

The right-hand side of \eqref{unref_id} coincides with the right-hand side of \eqref{class_id} at $\gamma = 1$, but the expansions are different.

